\documentclass[UTF8,a4paper,12pt]{ctexrep}

\usepackage[margin=2.35cm]{geometry}
\usepackage{xcolor}
\usepackage{hyperref}
\usepackage{booktabs}
\usepackage{longtable}
\usepackage{tabularx}
\usepackage{array}
\usepackage{enumitem}
\usepackage{listings}
\usepackage{amsmath}
\usepackage{tikz}
\usetikzlibrary{arrows.meta,positioning,fit,shapes.geometric}

\hypersetup{
  colorlinks=true,
  linkcolor=blue!55!black,
  urlcolor=blue!55!black,
  citecolor=blue!55!black
}

\definecolor{codebg}{RGB}{248,248,248}
\definecolor{notecolor}{RGB}{30,80,130}
\definecolor{warncolor}{RGB}{150,80,20}

\lstdefinestyle{codestyle}{
  basicstyle=\ttfamily\small,
  backgroundcolor=\color{codebg},
  frame=single,
  framerule=0.2pt,
  rulecolor=\color{black!25},
  breaklines=true,
  columns=fullflexible,
  keepspaces=true,
  showstringspaces=false
}
\lstset{style=codestyle}

\newcommand{\code}[1]{\texttt{\detokenize{#1}}}
\newcommand{\file}[1]{\texttt{\detokenize{#1}}}
\newcommand{\topic}[1]{\texttt{\detokenize{#1}}}
\newcommand{\rosnode}[1]{\texttt{\detokenize{#1}}}
\newcommand{\param}[1]{\texttt{\detokenize{#1}}}
\newcommand{\msg}[1]{\texttt{\detokenize{#1}}}
\newcommand{\pkg}[1]{\texttt{\detokenize{#1}}}
\newcommand{\actionname}[1]{\texttt{\detokenize{#1}}}

\newcolumntype{Y}{>{\raggedright\arraybackslash}X}

\newenvironment{notebox}
{\begin{center}\begin{minipage}{0.92\linewidth}\color{notecolor}\hrule\vspace{0.55em}}
{\vspace{0.35em}\hrule\end{minipage}\end{center}}

\newenvironment{warnbox}
{\begin{center}\begin{minipage}{0.92\linewidth}\color{warncolor}\hrule\vspace{0.55em}}
{\vspace{0.35em}\hrule\end{minipage}\end{center}}

\setlist[itemize]{leftmargin=2em,itemsep=0.22em}
\setlist[enumerate]{leftmargin=2em,itemsep=0.22em}
\renewcommand{\arraystretch}{1.22}

\title{Go2 机器狗 Waypoint 避障算法教程\\
\large 基于 ROS 2 Bridge、自研路点跟踪与 Unitree 避障接口的当前项目说明}
\author{RSSI/IMU Go2 项目}
\date{2026 年 6 月}

\begin{document}
\maketitle
\tableofcontents

\chapter{总体视角：当前避障算法到底是什么}

\section{一句话概括}

当前项目的避障算法不是完整 Nav2 导航栈，也不是让 Nav2 官方 waypoint follower 直接控制机器人。当前采用的是一个更适合 RSSI/IMU 实验采集的混合架构：

\[
\begin{aligned}
\text{waypoint 路线}
&\rightarrow \text{自研 odom 路点跟踪器}
\rightarrow \topic{/cmd_vel} \\
&\rightarrow \text{Go2 ROS 2 bridge}
\rightarrow \text{Unitree ObstaclesAvoidClient} \\
&\rightarrow \text{Go2 机载避障执行}
\end{aligned}
\]

这套设计的核心思想是：上层只规定 ``走哪条实验路线''，中间层根据 odom 生成速度命令，底层由宇树官方避障接口负责近距离防撞和绕障。这样可以避免在 RSSI/IMU 数据采集阶段过早陷入完整 Nav2 costmap、footprint、DWB critic、obstacle range 等大量细节调参。

\begin{notebox}
\textbf{当前代码主入口：}
\begin{lstlisting}
ros2_ws/src/go2_nav2_bridge/launch/go2_waypoint_obstacle_avoid.launch.py
\end{lstlisting}

它默认启动 \rosnode{go2_nav2_bridge} 和 \rosnode{go2_waypoint_follower}，并强制 bridge 使用 Unitree \code{obstacles_avoid} 运动后端。
\end{notebox}

\section{为什么不是“完整 Nav2 避障”}

完整 Nav2 避障通常是：

\[
\begin{aligned}
\text{目标点}
&\rightarrow \rosnode{planner_server}
\rightarrow \text{global path} \\
&\rightarrow \rosnode{controller_server}
\rightarrow \text{DWB/RPP}
\rightarrow \topic{/cmd_vel}
\end{aligned}
\]

这要求 Nav2 自己理解环境障碍，所以必须认真配置：

\begin{itemize}
  \item robot footprint 或 robot radius；
  \item local/global costmap；
  \item VoxelLayer 或 ObstacleLayer；
  \item obstacle range、raytrace range；
  \item inflation radius；
  \item controller 插件，如 DWB 或 Regulated Pure Pursuit；
  \item TF、odom、点云坐标系、速度/加速度限制。
\end{itemize}

而当前项目的实验目标首先是让 Go2 按可复现路线采集 RSSI/IMU 数据。对这个阶段来说，最重要的是路线可控、速度温和、日志可解释、出问题时容易定位。因此当前没有把 Nav2 controller/costmap 作为主避障链路。

\section{当前算法总图}

\begin{figure}[!htbp]
\centering
\begin{tikzpicture}[
  node distance=0.55cm and 0.8cm,
  main/.style={draw,rounded corners=2pt,align=center,text width=7.6cm,minimum height=0.82cm,fill=blue!5},
  exec/.style={draw,rounded corners=2pt,align=center,text width=7.6cm,minimum height=0.82cm,fill=orange!10},
  side/.style={draw,rounded corners=2pt,align=left,text width=4.6cm,minimum height=0.72cm,fill=gray!8},
  arr/.style={-{Latex[length=2.2mm]},thick}
]
  \node[main] (wp) {1. 路线输入：\code{relative_waypoints}\\L 型/多段实验路线};
  \node[main,below=of wp] (follower) {2. 自研 waypoint follower\\根据 \topic{/odom} 和当前目标点生成 \topic{/cmd_vel}};
  \node[main,below=of follower] (bridge) {3. Go2 ROS 2 bridge\\速度限幅、命令 watchdog、点云新鲜度门控};
  \node[exec,below=of bridge] (sdk) {4. Unitree \code{ObstaclesAvoidClient}\\\code{UseRemoteCommandFromApi(true)} + \code{SwitchSet(true)} + \code{Move}};
  \node[exec,below=of sdk] (go2) {5. Go2 机器狗执行\\机载近距离防撞与绕障};

  \node[side,right=of follower] (odom) {\topic{/odom}, \topic{/tf}\\来自 Go2 SportModeState};
  \node[side,right=of bridge] (cloud) {\topic{/go2/pointcloud}\\新鲜度门控、诊断、RViz};

  \draw[arr] (wp) -- (follower);
  \draw[arr] (follower) -- (bridge);
  \draw[arr] (bridge) -- (sdk);
  \draw[arr] (sdk) -- (go2);
  \draw[arr] (odom) -- (follower);
  \draw[arr] (cloud) -- (bridge);
\end{tikzpicture}
\caption{当前 Go2 waypoint 避障算法的数据流和控制流。}
\end{figure}

\section{模块分工}

\begin{table}[!htbp]
\centering
\small
\begin{tabularx}{\linewidth}{p{0.27\linewidth}p{0.14\linewidth}Y}
\toprule
模块 & 当前是否启用 & 作用 \\
\midrule
\rosnode{go2_waypoint_follower} & 是 & 读取相对 waypoint 和 \topic{/odom}，计算 \topic{/cmd_vel}。它负责路线跟踪，不负责点云避障。\\
\rosnode{go2_nav2_bridge} & 是 & 连接 ROS 2 与 Unitree SDK。发布 odom/IMU/点云，订阅 \topic{/cmd_vel} 并下发给 SDK。\\
Unitree \code{ObstaclesAvoidClient} & 是 & 当前近距离防撞和绕障的主要执行接口。\\
\topic{/go2/pointcloud} & 是 & bridge 发布点云，用于新鲜度门控、诊断、可视化；当前没有进入 Nav2 costmap 主链路。\\
Nav2 planner/controller/costmap & 参数保留 & 当前 waypoint launch 不启动。未来若需要完整 Nav2 避障，可以接入。\\
Nav2 官方 waypoint follower & 否 & 当前不使用，原因见第 \ref{chap:why-not-nav2-waypoint} 章。\\
\bottomrule
\end{tabularx}
\caption{当前 waypoint 避障算法的模块分工。}
\end{table}

\chapter{模块一：waypoint 路线如何表达和处理}

\section{relative\_waypoints 的格式}

当前路线通过 launch 参数 \param{relative_waypoints} 给出。它是一个数字列表，每三个数是一组：

\[
[x_1,\ y_1,\ yaw_1,\ x_2,\ y_2,\ yaw_2,\ \ldots]
\]

其中：

\begin{itemize}
  \item \(x_i\)：相对起点坐标系的前向距离，单位 m；
  \item \(y_i\)：相对起点坐标系的左向距离，单位 m；
  \item \(yaw_i\)：到达路点时希望的相对航向，单位 rad。
\end{itemize}

默认值是：

\begin{lstlisting}
relative_waypoints: [3.0, 0.0, 0.0]
\end{lstlisting}

含义是：从启动位置向前走 3 m，并保持启动时的航向。

\section{L 型路线示例}

比如要实现 ``先前进 3 m，到拐点，左转，再前进 2 m''，可以写成：

\begin{lstlisting}
ros2 launch go2_nav2_bridge go2_waypoint_obstacle_avoid.launch.py \
  network_interface:=eth0 \
  relative_waypoints:="[3.0, 0.0, 1.57, 3.0, 2.0, 1.57]"
\end{lstlisting}

解释如下：

\begin{itemize}
  \item 第一个路点 \([3.0, 0.0, 1.57]\)：相对起点向前 3 m，到达后朝左约 \(90^\circ\)。
  \item 第二个路点 \([3.0, 2.0, 1.57]\)：仍以起点坐标系为参考，目标位置是前方 3 m、左侧 2 m，航向保持左转后的方向。
\end{itemize}

\begin{warnbox}
当前所有 waypoint 都是相对于启动时的起点坐标系，不是相对于上一个 waypoint。也就是说第二个点不是 ``从第一个点再走 2 m''，而是 ``在起点坐标系下的绝对相对坐标''。
\end{warnbox}

\section{launch 如何解析 waypoint}

当前 waypoint launch 使用 \code{OpaqueFunction} 和 \code{ast.literal_eval} 解析 \param{relative_waypoints}。这样做比直接把字符串塞给参数更稳，因为它会检查：

\begin{itemize}
  \item 输入必须是 list 或 tuple；
  \item 长度不能为 0；
  \item 长度必须是 3 的倍数；
  \item 每个元素都能转成 float。
\end{itemize}

解析完成后，launch 把数值列表传给 \rosnode{go2_waypoint_follower}：

\begin{lstlisting}
Node(
  package="go2_nav2_bridge",
  executable="go2_waypoint_follower",
  parameters=[
    params_file,
    {"relative_waypoints": relative_waypoints},
  ],
)
\end{lstlisting}

\section{相对路点如何变成 odom 坐标}

\rosnode{go2_waypoint_follower} 第一次收到 \topic{/odom} 时，会记录当前位姿为起点：

\[
(x_0,\ y_0,\ \psi_0)
\]

之后每个相对路点 \((x_r,y_r,\psi_r)\) 转成 odom 坐标系下的目标：

\[
\begin{aligned}
x_g &= x_0 + \cos\psi_0 \cdot x_r - \sin\psi_0 \cdot y_r \\
y_g &= y_0 + \sin\psi_0 \cdot x_r + \cos\psi_0 \cdot y_r \\
\psi_g &= \mathrm{wrap}(\psi_0 + \psi_r)
\end{aligned}
\]

这一步的意义很大：实验人员不用知道场地中的绝对坐标，只要按机器人启动时的朝向设计路线即可。

\section{路点到达判定}

当前有两个容差：

\begin{itemize}
  \item \param{xy_tolerance_m}：位置误差小于该值，认为到达目标位置；
  \item \param{yaw_tolerance_rad}：航向误差小于该值，认为达到目标朝向。
\end{itemize}

只有位置和航向都满足要求，follower 才会切换到下一个 waypoint。全部 waypoint 完成后，节点发布零速度。

\chapter{模块二：自研 waypoint follower 如何算速度}

\section{输入输出}

\begin{longtable}{p{0.24\linewidth}p{0.28\linewidth}p{0.38\linewidth}}
\toprule
方向 & 名称 & 作用 \\
\midrule
输入 & \topic{/odom} & 当前机器人位置、姿态、速度。\\
输入 & \param{relative_waypoints} & 需要跟踪的路线。\\
输出 & \topic{/cmd_vel} & 发给 bridge 的速度命令。\\
参数 & \param{k_xy}, \param{k_yaw} & 位置和角度比例控制增益。\\
参数 & \param{max_vx_mps}, \param{max_vy_mps}, \param{max_wz_radps} & follower 侧速度上限。\\
\bottomrule
\end{longtable}

\section{控制律}

对当前目标点，follower 计算：

\[
d = \sqrt{(x_g-x)^2 + (y_g-y)^2}
\]

\[
\theta_e = \mathrm{wrap}(\mathrm{atan2}(y_g-y,x_g-x)-\psi)
\]

其中 \(d\) 是距离误差，\(\theta_e\) 是机器人当前朝向和目标方向之间的夹角。

当前速度计算为：

\[
\begin{aligned}
v_x &= \mathrm{clamp}(k_{xy}\,d\,\max(\cos\theta_e,0), 0, v_{x,\max})\\
v_y &= \mathrm{clamp}(k_{xy}\,d\,\sin\theta_e, -v_{y,\max}, v_{y,\max})\\
\omega_z &= \mathrm{clamp}(k_{yaw}\,\theta_e, -\omega_{\max}, \omega_{\max})
\end{aligned}
\]

这套控制律有几个直观效果：

\begin{itemize}
  \item 目标越远，前进速度越大；
  \item 目标越偏左，横向速度和角速度越偏左；
  \item 如果目标在身后，\(\cos\theta_e\) 可能为负，此时前向速度被压到 0，避免倒着硬走；
  \item 速度最终被上限裁剪，保证实验阶段运动温和。
\end{itemize}

\section{接近路点时怎么处理}

如果位置已经进入 \param{xy_tolerance_m}，但 yaw 还没达到 \param{yaw_tolerance_rad}，follower 不再平移，只发布角速度：

\[
v_x = 0,\quad v_y = 0,\quad \omega_z = \mathrm{clamp}(k_{yaw}(\psi_g-\psi),-\omega_{\max},\omega_{\max})
\]

这样 L 型路线在拐点处会先调整方向，再进入下一段运动。

\section{它负责什么，不负责什么}

\begin{table}[!htbp]
\centering
\small
\begin{tabularx}{\linewidth}{YY}
\toprule
\rosnode{go2_waypoint_follower} 负责 & \rosnode{go2_waypoint_follower} 不负责 \\
\midrule
解析路点列表；记录起点；将相对路点转换到 odom；按顺序跟踪 waypoint；生成平滑、低速、可解释的 \topic{/cmd_vel}。 &
读取点云；建立 costmap；对障碍物做全局重规划；像 DWB 一样对多组候选速度评分；直接调用 Unitree SDK。\\
\bottomrule
\end{tabularx}
\caption{自研 waypoint follower 的职责边界。}
\end{table}

避障执行发生在后面的 bridge 和 Unitree \code{ObstaclesAvoidClient}。

\chapter{模块三：go2\_nav2\_bridge 的接口转换与安全门控}

\section{bridge 为什么必要}

ROS 2 控制器输出的是 \topic{/cmd_vel}，Nav2 和 RViz 需要的是 \topic{/odom}、TF、点云等标准话题。但 Go2 侧原生接口来自 Unitree SDK/DDS，例如：

\begin{itemize}
  \item \code{rt/sportmodestate}：运动状态；
  \item \code{rt/utlidar/cloud}：点云；
  \item \code{SportClient::Move}：普通速度控制；
  \item \code{ObstaclesAvoidClient::Move}：带官方避障能力的速度控制。
\end{itemize}

\rosnode{go2_nav2_bridge} 是这两套接口之间的转换层。

\section{bridge 发布的 ROS 2 话题}

\begin{table}[!htbp]
\centering
\small
\begin{tabularx}{\linewidth}{p{0.23\linewidth}p{0.23\linewidth}Y}
\toprule
话题 & 类型 & 来源与作用 \\
\midrule
\topic{/odom} & \code{Odometry} & 来自 Go2 位置、速度和姿态。follower 依赖它闭环。\\
\topic{/tf} & TF & 发布 \code{odom -> base_link}。\\
\topic{/go2/imu} & \code{Imu} & 来自 Go2 IMU，供 RSSI/IMU 融合使用。\\
\topic{/go2/pointcloud} & \code{PointCloud2} & Unitree 点云经坐标修正和自车点过滤后发布。\\
\topic{/go2/path_s} & \code{Float32} & 相邻位置增量累加得到的路径长度估计。\\
\bottomrule
\end{tabularx}
\caption{bridge 发布的主要 ROS 2 话题。}
\end{table}

\section{速度命令如何进入 Unitree SDK}

follower 发布 \topic{/cmd_vel}，bridge 的 \code{onTwist()} 回调会调用 \code{sendVelocity(vx,vy,wz)}。这个函数做四件事：

\begin{enumerate}
  \item 把 NaN 或无穷值变成 0；
  \item 根据 \param{max_vx_mps}、\param{max_vy_mps}、\param{max_wz_radps} 限幅；
  \item 如果开启点云新鲜度门控，则检查 \topic{/go2/pointcloud} 是否超时；
  \item 根据 \param{motion_client} 选择 \code{ObstaclesAvoidClient::Move} 或 \code{SportClient::Move}。
\end{enumerate}

当前 waypoint launch 强制使用：

\begin{lstlisting}
motion_client: obstacles_avoid
obstacle_avoid_switch_on_start: true
obstacle_avoid_remote_api_on_start: true
\end{lstlisting}

因此最终运动调用是：

\begin{lstlisting}
obstacle_client_->Move(vx, vy, wz)
\end{lstlisting}

\section{安全门控}

\begin{longtable}{p{0.28\linewidth}p{0.62\linewidth}}
\toprule
门控 & 作用 \\
\midrule
速度限幅 & 防止上游 follower 或其他节点发出过大的速度命令。\\
非有限值处理 & 避免 NaN/Inf 进入 SDK。\\
点云新鲜度检查 & \param{require_fresh_point_cloud_for_motion=true} 时，非零速度必须建立在最近有点云更新的基础上。\\
命令 watchdog & 如果超过 \param{cmd_timeout_sec} 没收到新 \topic{/cmd_vel}，bridge 会发停止。\\
析构停止 & bridge 退出时调用 \code{sendStop()}，尽量保证机器人不残留速度。\\
\bottomrule
\end{longtable}

\section{点云处理}

当前 bridge 对 Unitree 点云做三步处理：

\begin{enumerate}
  \item 转换 DDS 点云到 ROS 2 \msg{PointCloud2}；
  \item 根据 \param{point_cloud_xyz_offset} 和 \param{point_cloud_rpy_offset} 修正坐标；
  \item 根据 \param{self_filter_boxes} 把自车附近点设为 NaN。
\end{enumerate}

这些处理对两件事有用：

\begin{itemize}
  \item 让 RViz 和诊断脚本看到更接近 \code{base_link} 坐标系的点云；
  \item 让点云新鲜度门控和未来 Nav2 costmap 接入更可靠。
\end{itemize}

但当前主避障链路里，\topic{/go2/pointcloud} 没有进入 Nav2 VoxelLayer/DWB。它更多是安全门控、诊断和未来扩展的基础。

\chapter{模块四：Unitree ObstaclesAvoidClient 的角色}

\section{为什么使用 obstacles\_avoid}

Go2 可以用普通运动接口发送速度，也可以用官方避障运动接口发送速度。当前 waypoint 避障算法选择后者，因为实验阶段不希望上层自己承担全部近距离避障细节。

启动时 bridge 会做：

\begin{lstlisting}
obstacle_client_->UseRemoteCommandFromApi(true);
obstacle_client_->SwitchSet(true);
\end{lstlisting}

含义是：

\begin{itemize}
  \item 允许 API 远程命令控制；
  \item 打开官方避障开关；
  \item 之后的速度命令通过 \code{ObstaclesAvoidClient::Move(vx,vy,wz)} 发送。
\end{itemize}

\section{它和上层 follower 的关系}

follower 输出的是 ``期望速度''，不是最终真实运动轨迹。Unitree \code{obstacles_avoid} 在执行时可能因为近距离障碍调整实际运动，甚至拒绝某些危险速度。

因此当前系统可以理解为两层：

\begin{itemize}
  \item 上层路线层：我希望去哪个 waypoint；
  \item 下层安全执行层：在不撞障碍的前提下尽量执行这个速度。
\end{itemize}

\section{它的优点和限制}

\begin{longtable}{p{0.22\linewidth}p{0.68\linewidth}}
\toprule
类型 & 说明 \\
\midrule
优点 & 避免在早期实验中大量调 Nav2 costmap/footprint/obstacle range。\\
优点 & 直接使用 Go2 官方运动能力，和机载感知/控制结合更紧。\\
优点 & 上层路线控制更简单，RSSI/IMU 数据采集更容易复现。\\
限制 & 官方避障是黑盒，无法像 Nav2 DWB critic 那样解释每个障碍点如何影响速度评分。\\
限制 & 如果官方避障改变了实际轨迹，odom 轨迹会偏离理想 waypoint 路线。\\
限制 & 它不是全局规划器。如果通道完全堵住，它不一定能规划绕远路。\\
\bottomrule
\end{longtable}

\chapter{为什么不采用 Nav2 官方 waypoint follower}
\label{chap:why-not-nav2-waypoint}

\section{Nav2 waypoint follower 实际做什么}

Nav2 的 \pkg{nav2_waypoint_follower} 名字里有 follower，但它不是底层速度控制器。它更像一个任务调度器：

\[
\begin{aligned}
\text{waypoint 列表}
&\rightarrow \pkg{nav2_waypoint_follower}
\rightarrow \actionname{NavigateToPose} \\
&\rightarrow \text{完整 Nav2 导航链路}
\end{aligned}
\]

真正把目标点变成 \topic{/cmd_vel} 的不是 waypoint follower，而是 Nav2 的 planner、controller、costmap、行为树等模块。

\section{如果直接用 Nav2 waypoint follower，会发生什么}

链路会变成：

\[
\begin{aligned}
\pkg{nav2_waypoint_follower}
&\rightarrow \rosnode{bt_navigator}
\rightarrow \rosnode{planner_server} \\
&\rightarrow \rosnode{controller_server}
\rightarrow \topic{/cmd_vel} \\
&\rightarrow \rosnode{go2_nav2_bridge}
\rightarrow \text{Go2 SDK}
\end{aligned}
\]

这时 Nav2 必须自己判断障碍和可通行空间，所以仍然绕不开：

\begin{itemize}
  \item footprint 是否准确；
  \item 点云是否正确落在 \code{base_link}/\code{odom}；
  \item obstacle range 和 raytrace range；
  \item inflation radius；
  \item DWB/RPP 控制器参数；
  \item recovery 行为是否适合机器狗。
\end{itemize}

这和当前项目想要的 ``尽量少调 Nav2 复杂避障参数，主要依靠 Go2 官方近距离避障'' 不一致。

\section{为什么当前选择自研 follower}

当前自研 follower 很小，只做一件事：根据 odom 追踪相对 waypoint，然后发速度。它的好处是：

\begin{itemize}
  \item 路线含义清楚：相对起点坐标；
  \item 行为可预测：按顺序到点、转向、继续；
  \item 调参少：主要调速度上限、比例增益、到点容差；
  \item 和 RSSI/IMU 实验匹配：路线复现比复杂全自主导航更重要；
  \item 不和 Unitree 官方避障抢同一层职责。
\end{itemize}

\section{这不是否定 Nav2}

当前方案不是说 Nav2 不好，而是分工不同。Nav2 更适合承担完整自主导航：建 costmap、规划路径、局部避障、恢复行为。当前 RSSI/IMU 项目阶段更需要一个可靠、简单、可解释的路线采集系统。

因此更准确的定位是：

\begin{quote}
当前项目使用 ROS 2/Nav2 兼容接口组织机器人状态和速度命令，但在线 waypoint 避障主链路采用自研 follower 加 Unitree 官方避障执行层，而不是 Nav2 官方 waypoint follower 加完整 Nav2 controller。
\end{quote}

\chapter{完整 Nav2 方案什么时候值得切换}

\section{适合切换的场景}

如果后续目标从 ``RSSI/IMU 路线采集'' 变成 ``复杂环境自主导航''，完整 Nav2 方案就更有价值。例如：

\begin{itemize}
  \item 需要绕过大障碍，而不是只做近距离防撞；
  \item 需要自动选择左右绕行；
  \item 需要动态重规划；
  \item 需要在论文中解释点云如何进入 costmap、如何影响速度评分；
  \item 需要复用 Nav2 的行为树、恢复行为、任务插件。
\end{itemize}

\section{完整 Nav2 方案的链路}

未来可以改成：

\[
\begin{aligned}
\text{waypoints}
&\rightarrow \pkg{nav2_waypoint_follower}
\rightarrow \rosnode{bt_navigator}
\rightarrow \rosnode{planner_server} \\
&\rightarrow \rosnode{controller_server}
\rightarrow \rosnode{velocity_smoother}
\rightarrow \rosnode{collision_monitor} \\
&\rightarrow \rosnode{go2_nav2_bridge}
\rightarrow \text{Unitree SDK}
\end{aligned}
\]

此时 \topic{/go2/pointcloud} 会进入 VoxelLayer：

\[
\begin{aligned}
\topic{/go2/pointcloud}
&\rightarrow \text{VoxelLayer}
\rightarrow \text{local costmap} \\
&\rightarrow \text{DWB/RPP}
\rightarrow \topic{/cmd_vel}
\end{aligned}
\]

但这也意味着必须投入时间调 costmap、footprint、obstacle range 和 controller 参数。

\section{两种方案对比}

\begin{longtable}{p{0.26\linewidth}p{0.31\linewidth}p{0.33\linewidth}}
\toprule
维度 & 当前方案 & 完整 Nav2 waypoint 方案 \\
\midrule
路线输入 & relative waypoint list & Nav2 waypoint action/list \\
速度生成 & 自研 follower & Nav2 controller，如 DWB/RPP \\
近距离避障 & Unitree \code{obstacles_avoid} & Nav2 costmap/controller + 可选 Unitree 兜底 \\
参数复杂度 & 低 & 高 \\
可解释性 & 控制逻辑清楚，避障执行黑盒 & Nav2 避障链路可解释，但配置复杂 \\
适合阶段 & RSSI/IMU 采集、路线复现 & 自主导航、复杂绕障、论文级 Nav2 分析 \\
\bottomrule
\end{longtable}

\chapter{如何运行和观察当前算法}

\section{编译}

\begin{lstlisting}
cd ~/桌面/RSSI/RSSI/code/ros2_ws
colcon build --packages-select go2_nav2_bridge
source install/setup.bash
\end{lstlisting}

\section{运行单段路线}

\begin{lstlisting}
ros2 launch go2_nav2_bridge go2_waypoint_obstacle_avoid.launch.py \
  network_interface:=eth0 \
  relative_waypoints:="[3.0, 0.0, 0.0]"
\end{lstlisting}

\section{运行 L 型路线}

\begin{lstlisting}
ros2 launch go2_nav2_bridge go2_waypoint_obstacle_avoid.launch.py \
  network_interface:=eth0 \
  relative_waypoints:="[3.0, 0.0, 1.57, 3.0, 2.0, 1.57]"
\end{lstlisting}

\section{观察话题}

\begin{lstlisting}
ros2 node list
ros2 topic hz /odom
ros2 topic echo /cmd_vel
ros2 topic hz /go2/pointcloud
ros2 topic echo /go2/path_s
ros2 run tf2_ros tf2_echo odom base_link
\end{lstlisting}

正常情况下应看到：

\begin{itemize}
  \item \rosnode{/go2_nav2_bridge}；
  \item \rosnode{/go2_waypoint_follower}；
  \item \topic{/odom} 持续更新；
  \item \topic{/cmd_vel} 在未到达目标前持续有速度；
  \item \topic{/go2/pointcloud} 按一定频率更新；
  \item bridge 日志中出现 \code{cmd_vel -> obstacles_avoid Move}。
\end{itemize}

\chapter{调参指南}

\section{路线跟踪参数}

\begin{table}[!htbp]
\centering
\small
\begin{tabularx}{\linewidth}{p{0.33\linewidth}p{0.15\linewidth}Y}
\toprule
参数 & 当前默认值 & 含义 \\
\midrule
\param{controller_frequency_hz} & 20.0 & follower 控制频率。\\
\param{max_vx_mps} & 0.30 & 最大前向速度。\\
\param{max_vy_mps} & 0.08 & 最大横向速度。\\
\param{max_wz_radps} & 0.45 & 最大角速度。\\
\param{min_vx_mps} & 0.10 & 非零前进时的最低速度。\\
\param{k_xy} & 0.45 & 距离误差到平移速度的比例增益。\\
\param{k_yaw} & 0.90 & 角度误差到角速度的比例增益。\\
\param{xy_tolerance_m} & 0.20 & 路点位置容差。\\
\param{yaw_tolerance_rad} & 0.30 & 路点航向容差。\\
\bottomrule
\end{tabularx}
\caption{waypoint follower 的主要路线跟踪参数。}
\end{table}

\section{安全参数}

\begin{table}[!htbp]
\centering
\small
\begin{tabularx}{\linewidth}{p{0.40\linewidth}Y}
\toprule
参数 & 调参建议 \\
\midrule
\param{cmd_timeout_sec} & 上游速度断流后的停止时间。过大不安全，过小可能误停。\\
\param{require_fresh_point_cloud_for_motion} & 真实机器人建议保持 \code{true}。\\
\param{point_cloud_motion_timeout_sec} & 点云超时阈值。若点云频率低，可略微增大，但不要无限放宽。\\
\param{self_filter_boxes} & 过滤机身附近点。过大可能删掉近距离墙面，过小可能保留自车点。\\
\param{max_vx_mps}, \param{max_vy_mps}, \param{max_wz_radps} & bridge 最终限速，应不大于 follower 侧或与其一致。\\
\bottomrule
\end{tabularx}
\caption{bridge 的主要安全参数。}
\end{table}

\section{问题排查}

\begin{longtable}{p{0.25\linewidth}p{0.65\linewidth}}
\toprule
现象 & 检查方向 \\
\midrule
机器人不动 & 检查 \topic{/odom} 是否有数据、\topic{/cmd_vel} 是否发布、bridge 是否使用 \code{obstacles_avoid}、SDK 网卡是否正确。\\
刚启动就停 & 多半是点云新鲜度门控触发，检查 \topic{/go2/pointcloud} 频率。\\
拐点处抖动 & 减小 \param{k_yaw} 或增大 \param{yaw_tolerance_rad}。\\
目标附近冲过头 & 减小 \param{min_vx_mps} 或 \param{k_xy}，增大 \param{xy_tolerance_m}。\\
路线偏离理想轨迹 & 可能是 odom 漂移、Unitree 避障主动绕开、地面打滑或起点 yaw 记录不稳。\\
\bottomrule
\end{longtable}

\end{document}
